% ============================================================
% Revisão Bibliográfica em Speech Analytics e Comunicação Clínica
% Resumo estruturado — Scoping Review (PRISMA-ScR / JBI)
% Rafael Baena Neto · Universidade São Francisco (USF)
% ============================================================

\documentclass[12pt, a4paper]{article}

\usepackage[utf8]{inputenc}
\usepackage[T1]{fontenc}
\usepackage[brazil]{babel}
\usepackage[
  top=2.5cm, bottom=2.5cm,
  left=3cm,  right=2.5cm
]{geometry}
\usepackage{setspace}
\usepackage{parskip}
\usepackage{titlesec}
\usepackage{booktabs}
\usepackage{array}
\usepackage{longtable}
\usepackage{hyperref}

\hypersetup{
  colorlinks = true,
  urlcolor   = blue,
  linkcolor  = black,
  citecolor  = black
}

\titleformat{\section}{\normalfont\bfseries}{\thesection.}{0.5em}{}
\titleformat{\subsection}{\normalfont\bfseries\itshape}{\thesubsection.}{0.5em}{}

% ============================================================
\begin{document}
% ============================================================

\begin{center}
  {\large\bfseries
    Speech Analytics e Comunicação Clínica:\\
    Uma Revisão Bibliográfica
  }\\[0.8em]
  Rafael Baena Neto\\
  Universidade São Francisco (USF)\\[0.4em]
  {\small Mestrado em Ciência de Dados em Saúde}
\end{center}

\bigskip
\hrule
\bigskip

% ============================================================
\begin{abstract}
\setlength{\parskip}{0.6em}
\onehalfspacing

\textbf{Contexto:}
A comunicação profissional-paciente constitui uma fonte densa de informação
clínica que ainda permanece predominantemente não estruturada e subutilizada
analiticamente. Avanços recentes em Reconhecimento Automático de Fala (ASR),
Processamento de Linguagem Natural (NLP) e análise de biomarcadores vocais
têm ampliado a capacidade de extrair dados estruturados a partir de interações
clínicas na anamnese. Ainda assim, poucos estudos operacionalizam essas
interações como variáveis mensuráveis derivadas da fala, especialmente no que
diz respeito a marcadores comunicacionais e comportamentais que podem ser
utilizados para análise de dados e avaliação de resultados no cuidado do
paciente.

\textbf{Objetivo:}
Identificar, sintetizar e avaliar criticamente as evidências sobre a aplicação
de \textit{Speech Analytics} e NLP à comunicação clínica profissional-paciente,
mapeando abordagens metodológicas, marcadores identificados, desafios técnicos
e lacunas para pesquisa futura.

\textbf{Métodos:}
Revisão bibliográfica. O corpus de análise foi composto por 27 artigos incluídos
após triagem de 36 registros provenientes de três bases documentais, organizados
em quatro eixos temáticos: (A)~NLP clínico e extração de entidades biomédicas;
(B)~documentação clínica automatizada e ASR; (C)~comunicação médico-paciente;
(D)~biomarcadores vocais e análise multimodal.

\textbf{Resultados:}
Os estudos evidenciam avanços consolidados em NLP clínico --- transição para
\textit{transformers}, BERT lidera tarefas de Reconhecimento de Entidades
Nomeadas (NER) --- e em transcrição automatizada (WER $\approx 9\%$ alcançável
com Whisper/PyAnnote). A comunicação centrada no paciente demonstra impacto
mensurável em desfechos objetivos de saúde. Marcadores comunicacionais como
autorreparo, engajamento, suporte emocional e indicadores de barreira são
identificados como preditores de adesão. A análise multimodal
(acústica\,+\,semântica) permanece em estado inicial. A maioria das abordagens
foca em automação documental, negligenciando a interação clínica como objeto
de análise em si.

\textbf{Conclusões:}
Confirma-se uma lacuna consistente: poucos estudos operacionalizam interações
clínicas como variáveis mensuráveis derivadas da fala, especialmente marcadores
comunicacionais e comportamentais extraídos da anamnese. Esta revisão
fundamenta a necessidade de estudos que integrem \textit{Speech Analytics} à
análise estruturada da comunicação clínica como fonte de dados para
monitoramento de adesão e qualidade do cuidado.

\end{abstract}

\bigskip
\hrule
\bigskip

\noindent\textbf{Palavras-chave:}
\textit{Speech Analytics};
revisão bibliográfica;
comunicação clínica;
NLP;
ASR;
marcadores comunicacionais;
\textit{digital scribe};
anamnese.

% ============================================================
\section{Perguntas de Pesquisa}
% ============================================================

\begin{itemize}
  \item \textbf{Q1:} Quais abordagens de \textit{Speech Analytics} têm sido
        aplicadas à comunicação clínica nos últimos cinco anos?
  \item \textbf{Q2:} Quais métodos de NLP e ASR são utilizados para transcrição
        e análise automatizada de consultas clínicas?
  \item \textbf{Q3:} Quais marcadores comunicacionais e comportamentais foram
        identificados ou propostos na interação profissional-paciente?
  \item \textbf{Q4:} Quais são os desafios técnicos, éticos e metodológicos
        para a análise automatizada da comunicação clínica?
\end{itemize}

% ============================================================
\section{Critérios de Inclusão e Exclusão}
% ============================================================

\begin{longtable}{p{1.8cm} p{10.5cm}}
\toprule
\textbf{Critério} & \textbf{Descrição} \\
\midrule
\endhead
IC1 & Estudo aplica \textit{Speech Analytics}, NLP ou ASR a contexto clínico \\[0.4em]
IC2 & Texto completo disponível em inglês, português ou espanhol \\[0.4em]
IC3 & Analisa comunicação profissional-paciente ou extração de informação de
      consultas clínicas \\[0.4em]
IC4 & Descreve métodos, resultados ou marcadores aplicáveis à interação clínica \\[0.4em]
IC5 & Publicação no período de janeiro de 2014 a 2026 \\
\midrule
EC1 & Publicações anteriores a 2014 \\[0.4em]
EC2 & Resumos estendidos, tutoriais ou \textit{white papers} sem dados empíricos \\[0.4em]
EC3 & Foco exclusivo em diagnóstico por voz de patologias físicas sem análise
      da interação comunicacional \\[0.4em]
EC4 & Estudos sobre populações não clínicas sem transferibilidade metodológica
      explicitada \\
\bottomrule
\end{longtable}

% ============================================================
% Referências
% ============================================================
\bigskip
\begin{thebibliography}{9}

\bibitem{page2021prisma}
PAGE, M.~J. et al.
The PRISMA~2020 statement: an updated guideline for reporting systematic
reviews.
\textit{BMJ}, v.~372, p.~n71, 2021.

\bibitem{tricco2018prisma}
TRICCO, A.~C. et al.
PRISMA Extension for Scoping Reviews (PRISMA-ScR): Checklist and
Explanation.
\textit{Annals of Internal Medicine}, v.~169, n.~7, p.~467--473, 2018.

\bibitem{peters2020jbi}
PETERS, M.~D.~J. et al.
Chapter 11: Scoping Reviews.
In: \textit{JBI Manual for Evidence Synthesis}. Adelaide: JBI, 2020.

\end{thebibliography}

\end{document}